% =========================================================================
% Figura: Interpretación geométrica del indicador de Hipervolumen (HV)
% Problema de minimización bidimensional
% =========================================================================

\begin{figure}[htbp]
\centering
\resizebox{0.58\linewidth}{!}{%
\begin{tikzpicture}[
    font=\sffamily,
    >=Stealth,
    x=1cm, y=1cm
]

% -------------------------------------------------------------------------
% 1. EJES COORDENADOS (Minimización)
% -------------------------------------------------------------------------
\draw[line width=1.1pt, draw=darkText, ->, -{Stealth[length=6pt, width=4pt]}]
    (0, 0) -- (6.8, 0) node[below=5pt, midway, font=\fontsize{9.5}{12}\selectfont, text=darkText] {\textbf{Objetivo $f_1$} (minimizar)};

\draw[line width=1.1pt, draw=darkText, ->, -{Stealth[length=6pt, width=4pt]}]
    (0, 0) -- (0, 5.7);
\node[font=\fontsize{9.5}{12}\selectfont, text=darkText, rotate=90]
    at (-0.65, 2.85) {\textbf{Objetivo $f_2$} (minimizar)};

% -------------------------------------------------------------------------
% 2. REGIÓN DEL HIPERVOLUMEN (Área sombreada escalonada)
% Soluciones A:
% a1=(1.0, 4.3), a2=(1.8, 3.3), a3=(2.8, 2.4), a4=(4.0, 1.4), a5=(5.1, 0.6)
% Punto de referencia: r = (5.8, 5.0)
% -------------------------------------------------------------------------
\filldraw[draw=figBlue!80, fill=figBlueTint, line width=1.0pt]
    (1.0, 5.0) -- (1.0, 4.3) --
    (1.8, 4.3) -- (1.8, 3.3) --
    (2.8, 3.3) -- (2.8, 2.4) --
    (4.0, 2.4) -- (4.0, 1.4) --
    (5.1, 1.4) -- (5.1, 0.6) --
    (5.8, 0.6) -- (5.8, 5.0) -- cycle;

% Líneas de proyección del punto de referencia
\draw[draw=grayText!70, dashed, line width=0.7pt] (5.8, 0) -- (5.8, 5.0);
\draw[draw=grayText!70, dashed, line width=0.7pt] (0, 5.0) -- (5.8, 5.0);
\node[below=3pt, font=\fontsize{8.5}{10}\selectfont, text=grayText] at (5.8, 0) {$r_1$};
\node[left=3pt, font=\fontsize{8.5}{10}\selectfont, text=grayText] at (0, 5.0) {$r_2$};

% Texto central en el hipervolumen
\node[text=figBlue!90!black, font=\fontsize{9.5}{12}\selectfont\bfseries, align=center] at (3.55, 4.15) {
    Región dominada\\
    \fontsize{8.5}{10.5}\selectfont Hipervolumen (\acrshort{hv})
};

% -------------------------------------------------------------------------
% 3. SOLUCIONES OBTENIDAS A (Círculos azules)
% -------------------------------------------------------------------------
\newcommand{\drawSol}[2]{%
    \filldraw[draw=figBlue!90!black, fill=figBlue, line width=0.9pt] (#1, #2) circle (0.12cm);%
}
\drawSol{1.0}{4.3}
\drawSol{1.8}{3.3}
\drawSol{2.8}{2.4}
\drawSol{4.0}{1.4}
\drawSol{5.1}{0.6}

% -------------------------------------------------------------------------
% 4. PUNTO DE REFERENCIA r (Círculo rojo/naranja)
% -------------------------------------------------------------------------
\filldraw[draw=figVermilion!90!black, fill=figVermilion, line width=0.9pt] (5.8, 5.0) circle (0.12cm);
\node[above=3pt, font=\fontsize{9}{11}\selectfont\bfseries, text=figVermilion!90!black]
    at (5.8, 5.0) {$\mathbf{r} = (r_1, r_2)$};

% -------------------------------------------------------------------------
% 5. LEYENDA (Ubicada en la esquina inferior izquierda)
% -------------------------------------------------------------------------
\begin{scope}[shift={(0.25, 0.25)}]
    \draw[draw=grayText!40, fill=white, rounded corners=3pt, line width=0.6pt]
        (0, 0) rectangle (2.50, 1.15);

    % Soluciones A
    \filldraw[draw=figBlue!90!black, fill=figBlue, line width=0.6pt] (0.22, 0.88) circle (0.07cm);
    \node[anchor=west, font=\fontsize{6.8}{8}\selectfont, text=darkText]
        at (0.38, 0.88) {Soluciones ($A$)};

    % Punto de referencia r
    \filldraw[draw=figVermilion!90!black, fill=figVermilion, line width=0.6pt] (0.22, 0.58) circle (0.07cm);
    \node[anchor=west, font=\fontsize{6.8}{8}\selectfont, text=darkText]
        at (0.38, 0.58) {Punto ref. ($\mathbf{r}$)};

    % Región sombreada
    \draw[draw=figBlue!80, fill=figBlueTint, line width=0.5pt] (0.14, 0.18) rectangle (0.30, 0.34);
    \node[anchor=west, font=\fontsize{6.8}{8}\selectfont, text=darkText]
        at (0.38, 0.26) {Área dominada};
\end{scope}

\end{tikzpicture}%
}
\caption{Interpretación geométrica del indicador de hipervolumen (\acrshort{hv}) en un espacio de dos objetivos a minimizar. El área sombreada delimita la región del espacio dominada por el conjunto de soluciones $A$ y acotada superiormente por el punto de referencia $\mathbf{r} = (r_1, r_2)$.}
\label{fig:indicador_hipervolumen}
\end{figure}
